\documentclass[10pt,a4paper]{amsart}
\usepackage{amsmath,amssymb,amsthm,graphicx}
% Retained minimal preamble for the Stacks Project chapter sites-cohomology.tex (Cohomology on Sites).
% Verbatim excerpt of the chapter preamble: Xy-pic setup (native preamble lines 14-18) and the native
% theorem declarations with the mathematical macros (native preamble lines 154-217). Copied mechanically;
% nothing added, removed or redefined. The wrapper class and the amsmath/amssymb/amsthm/graphicx base
% packages are supplied by the assembler, so the class/externaldocument part of the native preamble is omitted.
% For commutative diagrams we use Xy-pic
\usepackage[all]{xy}

% We use 2cell for 2-commutative diagrams.
\xyoption{2cell}
% Theorem environments.
%
\theoremstyle{plain}
\newtheorem{theorem}[subsection]{Theorem}
\newtheorem{proposition}[subsection]{Proposition}
\newtheorem{lemma}[subsection]{Lemma}

\theoremstyle{definition}
\newtheorem{definition}[subsection]{Definition}
\newtheorem{example}[subsection]{Example}
\newtheorem{exercise}[subsection]{Exercise}
\newtheorem{situation}[subsection]{Situation}

\theoremstyle{remark}
\newtheorem{remark}[subsection]{Remark}
\newtheorem{remarks}[subsection]{Remarks}

\numberwithin{equation}{subsection}

% Macros
%
\def\lim{\mathop{\mathrm{lim}}\nolimits}
\def\colim{\mathop{\mathrm{colim}}\nolimits}
\def\Spec{\mathop{\mathrm{Spec}}}
\def\Hom{\mathop{\mathrm{Hom}}\nolimits}
\def\Ext{\mathop{\mathrm{Ext}}\nolimits}
\def\SheafHom{\mathop{\mathcal{H}\!\mathit{om}}\nolimits}
\def\SheafExt{\mathop{\mathcal{E}\!\mathit{xt}}\nolimits}
\def\Sch{\mathit{Sch}}
\def\Mor{\mathop{\mathrm{Mor}}\nolimits}
\def\Ob{\mathop{\mathrm{Ob}}\nolimits}
\def\Sh{\mathop{\mathit{Sh}}\nolimits}
\def\NL{\mathop{N\!L}\nolimits}
\def\CH{\mathop{\mathrm{CH}}\nolimits}
\def\proetale{{pro\text{-}\acute{e}tale}}
\def\etale{{\acute{e}tale}}
\def\QCoh{\mathit{QCoh}}
\def\Ker{\mathop{\mathrm{Ker}}}
\def\Im{\mathop{\mathrm{Im}}}
\def\Coker{\mathop{\mathrm{Coker}}}
\def\Coim{\mathop{\mathrm{Coim}}}

% Boxtimes
%
\DeclareMathSymbol{\boxtimes}{\mathbin}{AMSa}{"02}

%
% Macros for moduli stacks/spaces
%
\def\QCohstack{\mathcal{QC}\!\mathit{oh}}
\def\Cohstack{\mathcal{C}\!\mathit{oh}}
\def\Spacesstack{\mathcal{S}\!\mathit{paces}}
\def\Quotfunctor{\mathrm{Quot}}
\def\Hilbfunctor{\mathrm{Hilb}}
\def\Curvesstack{\mathcal{C}\!\mathit{urves}}
\def\Polarizedstack{\mathcal{P}\!\mathit{olarized}}
\def\Complexesstack{\mathcal{C}\!\mathit{omplexes}}
% \Pic is the operator that assigns to X its picard group, usage \Pic(X)
% \Picardstack_{X/B} denotes the Picard stack of X over B
% \Picardfunctor_{X/B} denotes the Picard functor of X over B
\def\Pic{\mathop{\mathrm{Pic}}\nolimits}
\def\Picardstack{\mathcal{P}\!\mathit{ic}}
\def\Picardfunctor{\mathrm{Pic}}
\def\Deformationcategory{\mathcal{D}\!\mathit{ef}}

% --- appended 2026-10-12 (Chief): the project own definitions of the reference/slogan/history
% display environments, copied mechanically from the shared Stacks preamble lines 8-12
% (\usepackage{verbatim} plus the three \newenvironment lines, which the chapter preamble keeps
% verbatim from the project preamble). Nothing added, removed or redefined: these environments
% carry no mathematical content and the project itself typesets them as comments. Without these
% three lines a counted statement whose own text uses them fails with Environment slogan undefined.
\usepackage{verbatim}
\newenvironment{reference}{\comment}{\endcomment}
\newenvironment{slogan}{\comment}{\endcomment}
\newenvironment{history}{\comment}{\endcomment}
\usepackage{cleveref}
\crefname{originalreference}{original reference}{original references}
\Crefname{originalreference}{Original reference}{Original references}
\makeatletter
\newlabel{lemma-last-one-flat}{{0653}{source}{}{sourcepointer-1}{}}
\newlabel{lemma-last-one-flat@cref}{{[originalreference][0][]0653}{[1][0][]source}{}{}{}}
\makeatother
\crefname{originalsection}{original section}{original sections}
\Crefname{originalsection}{Original section}{Original sections}
\makeatletter
\newlabel{section-flat}{{00H9}{source}{}{sourcepointer-2}{}}
\newlabel{section-flat@cref}{{[originalsection][0][]00H9}{[1][0][]source}{}{}{}}
\makeatother
\crefname{originallemma}{original lemma}{original lemmas}
\Crefname{originallemma}{Original lemma}{Original lemmas}
\makeatletter
\newlabel{sites-modules-lemma-flat-ses}{{03EY}{source}{}{sourcepointer-3}{}}
\newlabel{sites-modules-lemma-flat-ses@cref}{{[originallemma][0][]03EY}{[1][0][]source}{}{}{}}
\makeatother
\title{Tor amplitude in a bounded interval is equivalent to representability by a bounded complex of flat modules}
\author{The Stacks Project Authors}
\date{stacks-project:sites-cohomology.tex:local-snapshot-20261011; GNU FDL 1.2 or later}
\begin{document}
\maketitle
\noindent\textit{Editorial note.} All delivered mathematical text is the author's own, copied from the pinned Stacks Project chapter file sites-cohomology.tex; no statement or proof step was rewritten, reordered or removed.\par\medskip
\noindent\textit{Editorial note.} The counted proof is the author's complete argument: the proof is written out in full: the first implication, the inductive truncation above b with the flatness of the kernel, and the truncation below a with the last-one-flat criterion.\par\medskip
\noindent\textit{Editorial note.} The label lemma-tor-amplitude names the counted statement, and the proof environment immediately following it in the source is delivered verbatim as its proof.\par\medskip
\noindent\textit{Editorial note.} The counted statement is carried by the author's own lemma environment in the chapter source.\par\medskip
\noindent\textit{Editorial note.} This card belongs to the priority queue of the D-72 twin ruling (the sites-side statement is the designated counting side, now counted here for the first time). The statement and proof are copied from the pinned chapter file; the three cited inputs are kept as external pointers with their canonical Stacks tags.\par\medskip
\noindent\textit{Allowed standard prerequisites (editorial).}\begin{itemize}
\item Existence of bounded above complexes of flat modules. Same chapter, the section on flatness providing a bounded above complex of flat O-modules representing a given object. Supplies the representative K\textasciicircum{}bullet used in the second implication.
\item Flatness of the kernel of a surjection of flat modules. Modules on Sites, the lemma stating that the kernel of a surjection between flat modules is flat. Supplies flatness of the kernel when truncating the representative.
\item The flatness criterion for the last term of a bounded above complex. Same chapter, the lemma concluding flatness of the degree-a cokernel from the vanishing of the tensor cohomology below a. Supplies flatness of E\textasciicircum{}a after truncation from below.
\end{itemize}\medskip
\ifdefined\hypertarget\hypertarget{sourcepointer-1}{}\fi
\noindent\textit{Source reference.} Original reference 0653: sources/stacks-sites-cohomology/pinned-original-source.tex, source lines 12403--12403. Prerequisite pointer invoked at proof line(s) 12467; it names a standard fact of the same chapter that the counted argument uses, and every step of the counted proof built on it is retained verbatim in the delivered body.\par\medskip
\ifdefined\hypertarget\hypertarget{sourcepointer-2}{}\fi
\noindent\textit{Source reference.} Original section 00H9: sources/stacks-sites-cohomology/pinned-original-source.tex, source lines 2961--2961. Prerequisite pointer invoked at proof line(s) 12453; it names a standard fact of the same chapter that the counted argument uses, and every step of the counted proof built on it is retained verbatim in the delivered body.\par\medskip
\ifdefined\hypertarget\hypertarget{sourcepointer-3}{}\fi
\noindent\textit{Source reference.} Original lemma 03EY: sources/stacks-chow/tags.txt, source lines 4473--4473. Prerequisite pointer invoked at proof line(s) 12458; it names a standard fact of a sibling chapter that the counted argument uses, and every step of the counted proof built on it is retained verbatim in the delivered body.\par\medskip
\begin{lemma}
\label{lemma-tor-amplitude}
Let $(\mathcal{C}, \mathcal{O})$ be a ringed site. Let $E$ be an object of
$D(\mathcal{O})$. Let $a, b \in \mathbf{Z}$ with $a \leq b$. The following
are equivalent
\begin{enumerate}
\item $E$ has tor-amplitude in $[a, b]$.
\item $E$ is represented by a complex
$\mathcal{E}^\bullet$ of flat $\mathcal{O}$-modules with
$\mathcal{E}^i = 0$ for $i \not \in [a, b]$.
\end{enumerate}
\end{lemma}

\begin{proof}
If (2) holds, then we may compute
$E \otimes_\mathcal{O}^\mathbf{L} \mathcal{F} =
\mathcal{E}^\bullet \otimes_\mathcal{O} \mathcal{F}$
and it is clear that (1) holds.

\medskip\noindent
Assume that (1) holds. We may represent $E$ by a bounded above complex
of flat $\mathcal{O}$-modules $\mathcal{K}^\bullet$, see
Section \ref{section-flat}.
Let $n$ be the largest integer such that $\mathcal{K}^n \not = 0$.
If $n > b$, then $\mathcal{K}^{n - 1} \to \mathcal{K}^n$ is surjective as
$H^n(\mathcal{K}^\bullet) = 0$. As $\mathcal{K}^n$ is flat we see that
$\Ker(\mathcal{K}^{n - 1} \to \mathcal{K}^n)$ is flat
(Modules on Sites, Lemma \ref{sites-modules-lemma-flat-ses}).
Hence we may replace $\mathcal{K}^\bullet$ by
$\tau_{\leq n - 1}\mathcal{K}^\bullet$. Thus, by induction on $n$, we
reduce to the case that $K^\bullet$ is a complex of flat
$\mathcal{O}$-modules with $\mathcal{K}^i = 0$ for $i > b$.

\medskip\noindent
Set $\mathcal{E}^\bullet = \tau_{\geq a}\mathcal{K}^\bullet$.
Everything is clear except that $\mathcal{E}^a$ is flat
which follows immediately from Lemma \ref{lemma-last-one-flat}
and the definitions.
\end{proof}

\end{document}
